\section{结论}\label{sec:conclusion}

我们在统一的评估协议下用仿真研究了一个问题：在现实噪声下 Kerr 微梳储备池表现如何、受什么限制、和软件相比处在什么位置。

第一，编码很重要。给输入光梳的每根梳齿加上随机复权重再送入微环，得到输出由输入决定、而不是由起振噪声决定的相干种子区。
在相同的物理条件和选点预算下，这种谱相位掩模在 Mackey--Glass 预测上比整体调制的光梳好一到两个数量级、比我们实现的单泵浦编码好两到三个数量级，
在信道均衡上与整体调制的光梳相当，单环在 NARMA-10 上没有帮助。梳齿相位的随机性比逐梳齿的幅度更重要，所以一个调制器加两个谱整形器就够了。
线性腔对照表明，Mackey--Glass 上的好处主要来自 Kerr 动力学；信道均衡上所需的非线性由相互干涉的输入梳齿经探测提供，Kerr 项没有贡献。

第二，下限由噪声决定，其中约一半在光源。理想探测器下误差下限约一半来自失谐抖动、一半来自散粒噪声；换 Si$_3$N$_4$ 或加功率都降不了它，
读出侧的改法（加大下载端口耦合、放慢符号）能降一部分。把真实抖动序列交给线性读出层没有去掉抖动，重新选工作点也没有绕开它，
所以本文器件的激光本征线宽要在几千赫兹以下；$100~\mathrm{kHz}$ 的普通激光让最好的单环比用 $1.6~\mathrm{kHz}$ 激光时差 5 到 6 倍。记忆要靠外加的延迟线和读出抽头。

第三，在软件坐标里，找到的最好微环低于同规模软件。它相当于输入接约 7 位转换器的 31 节点 ESN，而在微环能工作的约 $1~\mathrm{GBd}$ 下，现成的转换器有约 9 位；
把微环 31 路读数的量化也算进去后，在 6 到 12 位的每一档，只在输入端量化的 ESN 都更好。
16 个环拼接读出能把误差再降约 19 倍，但输入噪声更低的大型软件储备池同样能追平。
这个比较对给软件设定的输入噪声很敏感，所以我们给出整条曲线，并建议后续工作都这样报告。

第四，信号和抖动看来作用在同一个量（有效失谐）上，我们考察的工作点都避不开抖动。用反相信号同时驱动同一微环的两个方向（我们初步把好处归于交叉相位耦合），
读出两个方向并优化掩模，误差达到 $7.3\times10^{-6}$，比最好的单向微环低 21 倍，比同样读出、同样优化的单向微环低 4.5 倍；与输入权重得到同样优化的 62 节点 ESN 相比约为 8 位，比之前多 1 位，但仍低于这一速率下现成转换器的精度。

因此在本文探索的设计空间内，Kerr 微梳储备池的理由不在计算精度上。它能否落在物理层（信号本来在光域、省掉数字化、配合光域加权的延迟和能耗）上，本文没有检验；
这需要一个具体的光域应用、一次实验，以及在光域完成加权、只用一个转换器的读出方案。
